\section{Sample variance}\label{sec:sample-var}

Let $X_1,\ldots,X_n$ be i.i.d. copies of a centered variance-one random variable $X$, let $\bar X=n^{-1}\sum_iX_i$, and set
\[
Q_n=\sum_{i=1}^n(X_i-\bar X)^2.
\]
For $r\ge1$, let $V$ be uniform on $S^{r-1}$. Recall the spherical transform
\[
\Psi_r(z):=\E_V e^{zV_1}.
\]
Define
\[
\mathfrak D_{r,\tau}(Q):=
\sup_{|t|\le\tau}
\left|e^{-t^2/2}\E\Psi_r(t\sqrt Q)-1\right|.
\]
This is the spherical Laplace observable associated with a squared radial projection; it is not a probability metric on the law of $Q$.

The exact sample-variance consequence of the quadratic-form theorem is the following.
\begin{corollary}[Exact sample-variance characterization under dominance]\label{cor:exact-sample-var}
Let $n\ge2$ and let $X_1,\ldots,X_n$ be i.i.d., centered, and variance one. If $M_X$ is finite near the origin,
\[
M_X(s)M_X(-s)\ge e^{s^2}
\]
there, and
\[
\sum_{i=1}^n(X_i-\bar X)^2\sim\chi^2_{n-1},
\]
then $X_i\sim\Normal(0,1)$.
\end{corollary}

The quantitative statement is the next theorem.
\begin{theorem}[Sample-variance stability under spherical Laplace discrepancy]\label{thm:sample-var}
Fix $n\ge3$, $\tau,K>0$, and $0<\rho'<\tau\sqrt{(n-1)/n}$. There are constants $C<\infty$ and $\varepsilon_0>0$, depending only on $(n,\tau,K,\rho')$, such that the following holds. Let $X_1,\ldots,X_n$ be i.i.d. copies of a centered variance-one $X$ with $\E e^{2\tau|X|}\le K$, and assume $R_X(s)\ge0$ for $|s|\le\tau$. If
\[
Q_n=\sum_{i=1}^n(X_i-\bar X)^2,
\qquad
\mathfrak D_{n-1,\tau}(Q_n)\le\varepsilon\le\varepsilon_0,
\]
then
\[
\dK(X,\Normal(0,1))
\le C\sqrt{\frac{\log\log(1/\varepsilon)}{\log(1/\varepsilon)}}.
\]
\end{theorem}

The discrepancy-to-defect transfer is as follows. Let $r=n-1$, let $P=I_n-n^{-1}\mathbf1\mathbf1^{\mathsf T}$, and choose $B\in\R^{r\times n}$ with $B^{\mathsf T}B=P$ and $BB^{\mathsf T}=I_r$. If $V$ is uniform on $S^{r-1}$, then every column $b_i$ of $B$ has norm
\[
q:=\|b_i\|=\sqrt{\frac{n-1}{n}},
\]
and $\langle b_i,V\rangle$ has the same distribution as $qV_1$. Put $a=\tau q$ and define
\[
D_X(s):=\frac12R_X(s)
 =\frac12\log\frac{M_X(s)M_X(-s)}{e^{s^2}}.
\]
Under dominance, $D_X\ge0$ on $[-\tau,\tau]$.

The exact spherical transform gives, for $|t|\le\tau$,
\begin{equation}\label{eq:spherical-discrepancy-identity}
e^{-t^2/2}\E\Psi_r(t\sqrt{Q_n})
=\E_V\exp\left(\sum_{i=1}^n\left[\log M_X(t\langle b_i,V\rangle)-\frac{t^2\langle b_i,V\rangle^2}{2}\right]\right).
\end{equation}
Jensen's inequality, exchangeability, and the symmetry of $V_1$ therefore imply
\begin{equation}\label{eq:integrated-D}
n\E D_X(aV_1)\le\log\bigl(1+\mathfrak D_{n-1,\tau}(Q_n)\bigr).
\end{equation}
Equation~\eqref{eq:integrated-D} connects the spherical observable to one-dimensional reflection control.

For the pointwise transfer, the density of $V_1$ is
\[
f_n(v)=
\frac{\Gamma((n-1)/2)}{\sqrt\pi\,\Gamma((n-2)/2)}
(1-v^2)^{(n-4)/2},
\qquad -1<v<1.
\]
Fix $0<\rho'<a$ and choose $0<h<a-\rho'$. On $[-(\rho'+h),\rho'+h]$ the scaled density of $aV_1$ is bounded below by
\[
m_h:=\frac1a\frac{\Gamma((n-1)/2)}{\sqrt\pi\,\Gamma((n-2)/2)}
\left(1-\left(\frac{\rho'+h}{a}\right)^2\right)^{\max(n-4,0)/2}.
\]
The exponential envelope implies that $D_X$ is Lipschitz there. One may take
\[
L_h:=\frac{K}{e\tau}+\rho'+h,
\]
because $M_X(s)\ge1$ and $|x|e^{\tau|x|}\le(e\tau)^{-1}e^{2\tau|x|}$. If $M_*:=\sup_{|s|\le\rho'}D_X(s)$, choose $s_0\in[-\rho',\rho']$ with $D_X(s_0)=M_*$. The function $D_X$ is even, nonnegative, and satisfies $D_X(0)=0$. Since it is $L_h$-Lipschitz on the enlarged interval,
\[
D_X(s)\ge\bigl(M_*-L_h|s-s_0|\bigr)_+.
\]
Moreover, $|s_0|\ge M_*/L_h$, so the analogous tent centered at $-s_0$ has disjoint interior from the other one. Integrating the two tents against the density lower bound gives
\[
2m_h I_{L_h,h}(M_*)
 \le \frac1n\log\bigl(1+\mathfrak D_{n-1,\tau}(Q_n)\bigr),
\]
where
\[
I_{L,h}(M):=
\begin{cases}
M^2/L,& M\le Lh,\\
2hM-Lh^2,& M>Lh.
\end{cases}
\]
When $M\le Lh$, the integral is the area of a triangle of base $2M/L$; when $M>Lh$, it is the trapezoid obtained after the tent reaches the edge of the available interval. Writing $G(y;L,h)=\sqrt{Ly}$ when $y\le Lh^2$ and $G(y;L,h)=y/(2h)+Lh/2$ otherwise, we obtain
\begin{equation}\label{eq:pointwise-D}
\sup_{|s|\le\rho'}D_X(s)
\le\inf_{0<h<a-\rho'}G\left(\frac{\log(1+\mathfrak D_{n-1,\tau}(Q_n))}{2nm_h};L_h,h\right).
\end{equation}
In particular, for fixed $n,\tau,K,\rho'$, the right-hand side is $O(\sqrt{\mathfrak D_{n-1,\tau}(Q_n)})$ as the discrepancy tends to zero. The strict interior restriction is essential when $n\ge5$, because the spherical density vanishes at the endpoints.

Applying \eqref{eq:pointwise-D} to the sample-variance hypothesis gives
\[
\Delta_{\rho'}(X)
\le C_1\sqrt{\varepsilon},
\]
for all sufficiently small $\varepsilon$, where
$C_1=C_1(n,\tau,K,\rho')$; here we used $R_X=2D_X$.
Set $\Delta=C_1\sqrt{\varepsilon}$. Then
\[
\log(1/\Delta)=\tfrac12\log(1/\varepsilon)-\log C_1
\asymp \log(1/\varepsilon),
\qquad
\log\log(1/\Delta)=\log\log(1/\varepsilon)+O(1).
\]
Theorem~\ref{thm:sharp-upper}, with radius $\rho'$ and the fixed exponential envelope, therefore yields the displayed rate in Theorem~\ref{thm:sample-var}. In the exact case, Corollary~\ref{cor:exact-sample-var} gives Gaussianity under the same dominance condition. For $n=2$, the exact corollary is covered by the rank-one case of Theorem~\ref{thm:exact-psd}. Indeed, $A=I_2-\frac12\mathbf1\mathbf1^{\mathsf T}$ has rank one and both diagonal entries are positive. Thus the discrepancy-to-defect argument is a quantitative extension of the exact spherical-Jensen mechanism.

Two questions seem particularly natural. First, can the exponential-moment assumption in the reflection theorem be weakened without changing the optimal modulus? Second, under what additional uniform-integrability hypothesis can the spherical Laplace discrepancy in the sample-variance application be replaced by a standard probability metric? It would also be interesting to develop a corresponding quantitative theory for general positive-semidefinite quadratic forms with non-identically distributed coordinates.
